% Owner: Vinh
\subsection{Q3: MNIST Digit Classification}

% Vấn đề giải quyết
\subsubsection{Vấn đề giải quyết}
Bài toán yêu cầu phân loại ảnh chữ số viết tay trong tập dữ liệu MNIST
vào một trong 10 lớp, tương ứng với các chữ số từ 0 đến 9. Mỗi ảnh là ảnh
xám kích thước $28 \times 28$, với giá trị mỗi điểm ảnh nằm trong khoảng
$[0, 1]$. Ảnh được duỗi thành một vector đầu vào $x \in \mathbb{R}^{784}$.
Nhãn của mỗi ảnh được biểu diễn dưới dạng vector one-hot
$y \in \{0, 1\}^{10}$, trong đó chỉ phần tử ứng với chữ số đúng bằng 1.

Mô hình cần trả về vector điểm số (logits) $z \in \mathbb{R}^{10}$, và lớp
dự đoán là $\hat{c} = \arg\max_k z_k$. Mục tiêu là đạt độ chính xác tối
thiểu 97\% trên tập validation.

Khác với bài toán hồi quy ở Q2, nơi đầu ra là một giá trị liên tục, Q3 là
bài toán phân loại đa lớp: đầu ra là một lớp rời rạc trong số 10 lớp. Sự
khác biệt này quyết định cách chọn hàm mất mát cũng như cách diễn giải đầu
ra của mô hình, được trình bày ở các phần sau.

% Kiến thức sử dụng
\subsubsection{Kiến thức sử dụng}

\textbf{Mạng nơ-ron nhiều lớp (MLP).} Mô hình gồm các lớp tuyến tính
$f(x) = Wx + b$ xen kẽ với hàm kích hoạt phi tuyến
$\mathrm{ReLU}(u) = \max(0, u)$. Hàm phi tuyến là bắt buộc: nếu bỏ ReLU,
hai lớp tuyến tính liên tiếp $W_2(W_1x + b_1) + b_2$ vẫn chỉ là một ánh xạ
tuyến tính, nên mô hình không mạnh hơn một lớp duy nhất.

\textbf{Logits và softmax.} Đầu ra $z \in \mathbb{R}^{10}$ là các điểm số
chưa chuẩn hoá. Hàm softmax biến chúng thành phân phối xác suất trên 10 lớp:
\begin{equation}
p_k = \mathrm{softmax}(z)_k = \frac{e^{z_k}}{\sum_{j=1}^{10} e^{z_j}}.
\end{equation}

\textbf{Hàm mất mát cross-entropy.} Với nhãn one-hot $y$, cross-entropy đo
độ lệch giữa phân phối dự đoán $p$ và phân phối thật $y$:
\begin{equation}
L_{CE} = -\sum_{k=1}^{10} y_k \log p_k = -\log p_c,
\end{equation}
với $c$ là lớp đúng. Cực tiểu $L_{CE}$ tương đương với cực đại hợp lý
(maximum likelihood) của nhãn đúng. Đạo hàm theo logits có dạng rất gọn:
\begin{equation}
\frac{\partial L_{CE}}{\partial z_k} = p_k - y_k.
\end{equation}

\textbf{So sánh với MSE.} Hàm MSE dùng ở Q2 cũng có thể áp dụng trực tiếp
lên logits và nhãn one-hot:
\begin{equation}
L_{MSE} = \frac{1}{10}\sum_{k=1}^{10} (z_k - y_k)^2, \qquad
\frac{\partial L_{MSE}}{\partial z_k} = \frac{2}{10}(z_k - y_k).
\end{equation}
Trong thực nghiệm, mô hình huấn luyện bằng MSE vẫn đạt độ chính xác cao,
nhưng hội tụ chậm hơn đáng kể so với cross-entropy. Có thể giải thích như
sau:
\begin{itemize}
  \item \emph{Vì sao MSE vẫn hoạt động:} lớp dự đoán là $\arg\max_k z_k$.
  MSE kéo logit của lớp đúng về 1 và các logit còn lại về 0, nên khi
  $L_{MSE}$ đủ nhỏ thì logit lớp đúng là lớn nhất và dự đoán đúng.
  \item \emph{Độ lớn gradient:} PyTorch tính trung bình MSE trên toàn bộ
  $B \times 10$ phần tử, còn cross-entropy chỉ trung bình trên $B$ mẫu. Hệ số
  $2/10$ khiến gradient của MSE nhỏ hơn khoảng 5 lần, tương đương huấn luyện
  với learning rate nhỏ hơn.
  \item \emph{Mục tiêu không phù hợp:} MSE coi bài toán là hồi quy về đúng
  các giá trị 0 và 1. Một dự đoán đúng và rất tự tin (ví dụ $z_c = 5$,
  $z_k = -3$ với $k \neq c$) vẫn bị phạt, dù với bài toán phân loại đó là
  kết quả tốt. Một phần quá trình học bị dành cho mục tiêu không cần thiết.
  \item \emph{Tín hiệu học:} gradient $p_k - y_k$ của cross-entropy lớn khi
  mô hình dự đoán sai một cách tự tin, và giảm dần về 0 khi lớp đúng đã
  được tách biệt rõ. Đây đúng là hành vi cần có cho bài toán phân loại.
\end{itemize}
Lưu ý: nhận định thường gặp ``MSE kết hợp softmax gây triệt tiêu gradient''
chỉ đúng khi MSE được tính trên xác suất $p$ (do đạo hàm softmax chứa thừa
số $p_k(1 - p_k)$). Ở đây MSE được tính trên logits $z$, nên nguyên nhân
chính là độ lớn gradient và mục tiêu không phù hợp nêu trên.

\textbf{Mini-batch SGD.} Tham số được cập nhật theo
$\theta \leftarrow \theta - \eta \nabla_\theta L$, với gradient tính trung
bình trên một batch nhỏ $B$ mẫu. Batch nhỏ cho nhiều lần cập nhật hơn trong
một epoch, đổi lại gradient nhiễu hơn.

% Implement, công thức
\subsubsection{Cài đặt}

\textbf{Kiến trúc.} Mô hình là MLP hai lớp
$784 \rightarrow 128 \rightarrow 10$:
\begin{align}
h &= \mathrm{ReLU}(W_1 x + b_1), & W_1 &\in \mathbb{R}^{128 \times 784},\\
z &= W_2 h + b_2, & W_2 &\in \mathbb{R}^{10 \times 128}.
\end{align}
Tổng số tham số là $784 \cdot 128 + 128 + 128 \cdot 10 + 10 = 101\,770$.
Mô hình trả về logits $z$; hàm \texttt{cross\_entropy} của PyTorch đã tự áp
dụng softmax bên trong nên không cần thêm softmax ở đầu ra.

\textbf{Huấn luyện.} Mô hình được huấn luyện bằng SGD với learning rate
$\eta = 0.01$, batch size $B = 5$, tối đa 50 epoch. Sau mỗi epoch, độ chính
xác trên tập kiểm định được đánh giá, và quá trình huấn luyện dừng sớm
(early stopping) khi đạt ngưỡng 97\%.

\textbf{Kết quả.} Bảng~\ref{tab:q3_loss} so sánh hai hàm mất mát với cùng
kiến trúc và siêu tham số. Dòng cuối tăng learning rate của MSE lên 5 lần để
kiểm tra giả thuyết về độ lớn gradient.
\begin{table}[htbp]
\caption{So sánh cross-entropy và MSE trên MNIST}
\label{tab:q3_loss}
\centering
\begin{tabular}{lccc}
\toprule
Hàm mất mát & $\eta$ & Số epoch đạt 97\% & Độ chính xác cuối \\
\midrule
Cross-entropy & 0.01 & 6& 0.972 (97,2\%)\\
MSE           & 0.01 & ~33& 0.9692 (96,92\%)\\
MSE           & 0.05 & ~8& 0.9674 (96,74\%)\\
\bottomrule
\end{tabular}
\end{table}
